\documentclass{article}

\PassOptionsToPackage{numbers,sort&compress}{natbib}

\usepackage[eandd, final]{neurips_2026}
\usepackage[utf8]{inputenc} % allow utf-8 input
\usepackage[T1]{fontenc}    % use 8-bit T1 fonts
\usepackage{hyperref}       % hyperlinks
\usepackage{url}            % simple URL typesetting
\usepackage{booktabs}       % professional-quality tables
\usepackage{amsfonts}       % blackboard math symbols
\usepackage{nicefrac}       % compact symbols for 1/2, etc.
\usepackage{microtype}      % microtypography
\usepackage{xcolor}         % colors
\usepackage{paralist}
\usepackage{amsmath}

\newcommand{\pending}[1]{\textcolor{red}{[#1]}} % Open items. Remove before submission.

\title{Empirical Benchmark of Steering Methods for LLM Bias Mitigation}

% The \author macro works with any number of authors. There are two commands
% used to separate the names and addresses of multiple authors: \And and \AND.

% Using \And between authors leaves it to LaTeX to determine where to break the
% lines. Using \AND forces a line break at that point. So, if LaTeX puts 3 of 4
% authors names on the first line, and the last on the second line, try using
% \AND instead of \And before the third author name.

\author{%
  \makebox[0.4\textwidth]{Chin Ling Xing, Lance*}\\
  A0309152L\\
  e1406571\\
  \And
  \makebox[0.4\textwidth]{Lim Yan Khai, Bennett*}\\
  A0291595N\\
  e1332641\\
  \AND
  \makebox[0.4\textwidth]{Pradyumna Shyama Prasad*}\\
  A0276991J\\
  e1136918\\
  \And
  \makebox[0.4\textwidth]{Ou Chuhao*}\\
  A0211621Y\\
  e0496099\\
}

\begin{document}


\maketitle
{\renewcommand\thefootnote{*}\footnotetext{These authors are listed in no particular order.}}


\begin{abstract}
 % We will use Neurips style for the report. Ensure the main report is 6 pages long.

\end{abstract}



\section{Introduction}

Large Language Models (LLMs) have rapidly evolved over the past few years, demonstrating remarkable performance across numerous text generation tasks such as translation, summarization and creative writing. However, as these models were trained on text written by humans, they inherently perpetuate and exhibit similar societal biases present in the training data (Hu et al., 2025). Mitigating these biases therefore becomes important to ensure that LLMs do not produce harmful outputs that marginalize people across demographic lines, especially when they are deployed in real-world scenarios. Traditional bias mitigation strategies often involve dataset filtering, Supervised Fine-Tuning (SFT), or Reinforcement Learning from Human Feedback (RLHF), but these would permanently alter the model’s weights and possibly impair its original knowledge base to enforce safety. On the other hand, steering methods present a promising alternative, where the LLM’s internal activations are modified so that the model can frame its replies more neutrally. As activation steering is done during inference time, it is also computationally cheaper as compared to the traditional methods discussed. In this work, we seek to compare the performance of four well-known activation steering techniques based on how well they mitigate biases in their responses.


Hence, our problem statement is:
\begin{quote}
Can steering methods reliably mitigate stereotypes and biases in LLMs without degrading the quality of their text generations? 
\end{quote}

\section{Related Work}

\section{Methodology}

\section{Experiments \& Results}

\section{Conclusion \& Future Work}

\begin{ack}
% This section can go with the references and not be in the main 6-pages.

The authors utilized ChatGPT (GPT-5.6 Sol) and Claude (Opus 5.5) to implement the steering methods and evaluation scripts described in this paper. All code generated by these tools was systematically reviewed and validated by the authors to ensure correctness and consistency.

These AI tools were also used during the preparation of this manuscript, and every AI-generated section was subsequently revised and fact-checked by the authors.
\end{ack}

\section*{References}
% References follow the acknowledgments. Any choice of citation style is acceptable as long as you are
% consistent. It is permissible to reduce the font size to \verb+small+ (9 point)
% when listing the references. References can use a bib file and include the citation using the \verb+\cite{key}+ to cite the source in-line.
% Note that the Reference section does not count towards the page limit.

{
\small

Cunningham, H., Ewart, A., Riggs, L., Huben, R., \& Sharkey, L.\ (2023). Sparse autoencoders find highly interpretable features in language models. {\it arXiv preprint arXiv:2309.08600}.

Dong, Q., Li, L., Dai, D., Zheng, C., Ma, J., Li, R., ... \& Sui, Z. (2024). A survey on in-context learning. {\it In Proceedings of the 2024 conference on empirical methods in natural language processing}, 1107--1128.

Dubois, Y., Galambosi, B., Liang, P., \& Hashimoto, T.\ B.\ (2024). Length-controlled AlpacaEval: A simple way to debias automatic evaluators. {\it arXiv preprint arXiv:2404.04475}.

Gemma Team (2025). Gemma 3 technical report. {\it arXiv preprint arXiv:2503.19786}.

Hu, T., Kyrychenko, Y., Rathje, S., Collier, N., Van Der Linden, S., \& Roozenbeek, J. (2025). Generative language models exhibit social identity biases. {\it Nature Computational Science}, 5(1), 65--75.

Kaneko, M., Bollegala, D., \& Baldwin, T.\ (2026). A multilingual social bias benchmark incorporating thinking processes. {\it Proceedings of the 64th Annual Meeting of the Association for Computational Linguistics (Volume 1: Long Papers)}, 47726--47741.

Liu, S., Ye, H., Xing, L., \& Zou, J.\ Y.\ (2024). In-context vectors: Making in-context learning more effective and controllable through latent space steering. {\it Proceedings of the 41st International Conference on Machine Learning}, 235, 32287--32307.

Parrish, A., Chen, A., Nangia, N., Padmakumar, V., Phang, J., Thompson, J., Htut, P.\ M., \& Bowman, S.\ R.\ (2022). BBQ: A hand-built bias benchmark for question answering. {\it Findings of the Association for Computational Linguistics: ACL 2022}, 2086--2105.

Rimsky, N., Gabrieli, N., Schulz, J., Tong, M., Hubinger, E., \& Turner, A.\ M.\ (2024). Steering Llama 2 via contrastive activation addition. {\it Proceedings of the 62nd Annual Meeting of the Association for Computational Linguistics (Volume 1: Long Papers)}, 15504--15522.

Smith, E.\ M., Hall, M., Kambadur, M., Presani, E., \& Williams, A.\ (2022). ``I'm sorry to hear that'': Finding new biases in language models with a holistic descriptor dataset. {\it Proceedings of the 2022 Conference on Empirical Methods in Natural Language Processing}, 9180--9211.

Wang, Y., Ma, X., Zhang, G., Ni, Y., Chandra, A., Guo, S., Ren, W., Arulraj, A., He, X., Jiang, Z., Li, T., Ku, M., Wang, K., Zhuang, A., Fan, R., Yue, X., \& Chen, W.\ (2024). MMLU-Pro: A more robust and challenging multi-task language understanding benchmark. {\it Advances in Neural Information Processing Systems}, 37.

Wu, Z., Arora, A., Geiger, A., Wang, Z., Huang, J., Jurafsky, D., Manning, C.\ D., \& Potts, C.\ (2025). AxBench: Steering LLMs? Even simple baselines outperform sparse autoencoders. {\it Proceedings of the 42nd International Conference on Machine Learning}.
}


%%%%%%%%%%%%%%%%%%%%%%%%%%%%%%%%%%%%%%%%%%%%%%%%%%%%%%%%%%%%

\appendix

\section{Technical appendices and supplementary material}
You can include any number of technical appendices and supplementary material pages after the main report.

Work with the assumption that the appendix will not be graded. It is there to provide additional evidence. The main report (6-page-long) should be independent of the appendix.

%%%%%%%%%%%%%%%%%%%%%%%%%%%%%%%%%%%%%%%%%%%%%%%%%%%%%%%%%%%%


% NOTE: Additional pages can be used only for references/appendix.

\iffalse
\section{General formatting instructions}
\label{gen_inst}

You can use compact item for listing itemized texts, like below
\begin{compactitem}
\item item 1
\item item 2
\item item 3
\end{compactitem}

Normal itemized list uses itemize

\begin{itemize}
\item item 1
\item item 2
\item item 3
\end{itemize}

Similarly, you can use compact enum vs numbered list 

\begin{compactenum}
\item item 1
\item item 2
\item item 3
\end{compactenum}

\begin{enumerate}
\item item 1
\item item 2
\item item 3
\end{enumerate}

\section{Headings: first level}
\label{headings}

\textbf{If you are simply following the template with no modifications, everything should work as expected.}

All headings should be lower case (except for first word and proper nouns),
flush left, and bold.


First-level headings should be in 12-point type.


\subsection{Headings: second level}


Second-level headings should be in 10-point type.


\subsubsection{Headings: third level}


Third-level headings should be in 10-point type.


\paragraph{Paragraphs}


There is also a \verb+\paragraph+ command available, which sets the heading in
bold, flush left, and inline with the text, with the heading followed by 1\,em
of space.


\section{Citations, figures, tables, references}
\label{others}


These instructions apply to everyone.


\subsection{Citations within the text}


The \verb+natbib+ package will be loaded for you by default.  Citations may be
author/year or numeric, as long as you maintain internal consistency.  

Of note is the command \verb+\citet+, which produces citations appropriate for
use in inline text.  For example,
\begin{verbatim}
   \citet{hasselmo} investigated\dots
\end{verbatim}
produces
\begin{quote}
  Hasselmo, et al.\ (1995) investigated\dots
\end{quote}


If \verb+natbib+ clashes with another package you load, you can add the optional
argument \verb+nonatbib+ when loading the style file:
\begin{verbatim}
   \usepackage[nonatbib]{neurips_2026}
\end{verbatim}


\subsection{Footnotes}


Footnotes should be used sparingly.  If you do require a footnote, indicate
footnotes with a number\footnote{Sample of the first footnote.} in the
text. Place the footnotes at the bottom of the page on which they appear.
Precede the footnote with a horizontal rule of 2~inches (12~picas).


Note that footnotes are properly typeset \emph{after} punctuation
marks.\footnote{As in this example.}


\subsection{Figures}


\begin{figure}
  \centering
  \fbox{\rule[-.5cm]{0cm}{4cm} \rule[-.5cm]{4cm}{0cm}}
  \caption{Sample figure caption. Explain what the figure shows and add a key take-away message to the caption.}
\end{figure}


All artwork must be neat, clean, and legible. Lines should be dark enough to be readable. The figure number and caption always appear after the
figure. Place one line space before the figure caption and one line space after
the figure. The figure caption should be lowercase (except for the first word and proper nouns); figures are numbered consecutively.


\subsection{Tables}


All tables must be centered, neat, clean, and legible.  The table number and
title always appear before the table.  See Table~\ref{sample-table}.


Place one line space before the table title, one line space after the
table title, and one line space after the table. The table title must
be lowercase (except for the first word and proper nouns); tables are
numbered consecutively.

We strongly suggest the use of the \verb+booktabs+ package, which allows for
typesetting high-quality, professional tables:

This package was used to typeset Table~\ref{sample-table}.


\begin{table}
  \caption{Sample table caption. Explain what the table shows and add a key take-away message to the caption.}
  \label{sample-table}
  \centering
  \begin{tabular}{lll}
    \toprule
    \multicolumn{2}{c}{Part}                   \\
    \cmidrule(r){1-2}
    Name     & Description     & Size ($\mu$m) \\
    \midrule
    Dendrite & Input terminal  & $\approx$100     \\
    Axon     & Output terminal & $\approx$10      \\
    Soma     & Cell body       & up to $10^6$  \\
    \bottomrule
  \end{tabular}
\end{table}

\subsection{Math}
Please use LaTeX (or AMSTeX) commands for unnumbered display math. 
Inline math can be like this: $a + b = c$.

Centered math this way:

$$a^2 + b^2 = c^2$$

Numbered equations use \verb+equation+ environment

\begin{equation}
a^x + b^y = c^z
\label{eq1}
\end{equation}

You can refer to the equation using Eq.~\ref{eq1}.
\fi

\end{document}